%==============================================================================
%  WEEKLY AI BRIEF -- issue #11, week ending 3 November 2026
%  Engine: tectonic
%==============================================================================
\documentclass[11pt,a4paper]{article}

%------------------------------ packages -------------------------------------
\usepackage[T1]{fontenc}
\usepackage[utf8]{inputenc}
\usepackage{lmodern}
\usepackage[margin=2.2cm,top=2.0cm,bottom=2.3cm]{geometry}
\usepackage{microtype}
\usepackage{graphicx}
\usepackage{xcolor}
\usepackage{titlesec}
\usepackage{enumitem}
\usepackage{booktabs}
\usepackage{tabularx}
\usepackage{fancyhdr}
\usepackage{lastpage}
\usepackage{tcolorbox}
\usepackage{xurl}
\usepackage{needspace}
\usepackage[hidelinks]{hyperref}

\tcbuselibrary{skins,breakable}

% Logos live in template/ and are shared by every issue -- never copy them into a
% dated folder. This path list resolves them from <YYYY-MM-DD>/ or from template/.
\graphicspath{{../../../template/}{./}{template/}}

%------------------------------ palette --------------------------------------
\definecolor{brandink}{HTML}{103868}    % KAUST navy -- headings, rules
\definecolor{brandaccent}{HTML}{0B7480} % KAUST teal (darkened to 5.6:1 on white)
\definecolor{brandsoft}{HTML}{EAF3F5}   % box fills
\definecolor{brandgrey}{HTML}{5A6672}   % meta text

\hypersetup{colorlinks=true,linkcolor=brandaccent,urlcolor=brandaccent,
            citecolor=brandaccent,pdfborder={0 0 0}}

%------------------------------ metadata -------------------------------------
% Issue number: one more than the highest \briefnumber in the existing dated folders.
% A REBUILD of a window that already has a folder keeps that folder's number -- the
% number identifies the window, not the run.
\newcommand{\briefnumber}{11}
% The window is seven days ending on \briefdate, INCLUSIVE -- so the span is 6 days
% wide (19--25, not 18--25). Check both endpoints against `date +%F` before shipping.
\newcommand{\briefweek}{28 October -- 3 November 2026}
\newcommand{\briefdate}{3 November 2026}            % date of publication
\newcommand{\briefauthor}{Ali Habibullah}          % prepared by
\newcommand{\briefaffil}{KAUST}                    % affiliation
\newcommand{\briefscope}{Papers, model releases, benchmarks and policy from the past seven days}

%------------------------------ section styling ------------------------------
\titleformat{\section}
  {\normalfont\sffamily\large\bfseries\color{brandink}}{\thesection.}{0.6em}{}
  [\vspace{-0.55em}{\color{brandaccent}\rule{\textwidth}{0.9pt}}]
\titlespacing*{\section}{0pt}{1.5ex plus .3ex}{0.9ex}

\titleformat{\subsection}[runin]
  {\normalfont\sffamily\bfseries\color{brandink}}{}{0pt}{}[\;---]
\titlespacing*{\subsection}{0pt}{1.25ex plus .2ex}{0.5em}

\setlist[itemize]{leftmargin=1.15em,itemsep=0.25em,topsep=0.35em,parsep=0pt}
\setlength{\parindent}{0pt}
\setlength{\parskip}{0.45em}
\setlength{\emergencystretch}{1.5em}

%------------------------------ header / footer ------------------------------
\pagestyle{fancy}
\fancyhf{}
\renewcommand{\headrulewidth}{0pt}
\renewcommand{\footrulewidth}{0.4pt}
\fancyfoot[L]{\footnotesize\color{brandgrey}Weekly AI Brief \#\briefnumber\ \textbullet\ \briefweek}
\fancyfoot[R]{\footnotesize\color{brandgrey}Page \thepage\ of \pageref{LastPage}}

%------------------------------ content macros -------------------------------
% Executive-summary box.  \tldr{ \item ... \item ... }
\newcommand{\tldr}[1]{%
  \begin{tcolorbox}[breakable,enhanced,colback=brandsoft,colframe=brandaccent,
      boxrule=0.8pt,left=10pt,right=10pt,top=8pt,bottom=8pt,arc=2pt,
      title={\sffamily\bfseries The week in five lines},fonttitle=\color{white}]
    \begin{itemize}[leftmargin=1.1em,itemsep=0.3em,topsep=0pt]#1\end{itemize}
  \end{tcolorbox}}

% One item of news.  \entry{Headline}{Who / venue}{URL}{source-number}
\newcommand{\entry}[4]{%
  \par\Needspace{7\baselineskip}\vspace{0.7em}%
  {\sffamily\bfseries\color{brandink}#1}\par
  {\footnotesize\color{brandgrey}#2 \textbullet\ \href{#3}{link} \textbullet\ \srcref{#4}}\par
  \vspace{0.15em}}

% "Why it matters" line -- one sentence, no hedging.
\newcommand{\why}[1]{\par{\small\textbf{Why it matters.} #1}\par}

% Numbered pointer into SOURCES.md / the source table.
\newcommand{\srcref}[1]{\textcolor{brandaccent}{[S#1]}}

% Numbers worth pulling out of a paragraph.
\newcommand{\stat}[2]{{\sffamily\bfseries\color{brandink}#1}~{\footnotesize\color{brandgrey}#2}}

%==============================================================================
\begin{document}

%------------------------------ title block ----------------------------------
\noindent
\begin{minipage}[c]{0.48\textwidth}
  \raggedright\includegraphics[height=1.4cm]{KAUST_Academy_Logo.png}
\end{minipage}\hfill
\begin{minipage}[c]{0.48\textwidth}
  \raggedleft\includegraphics[height=1.5cm]{KAUST_Logo.png}
\end{minipage}
\vspace{0.7em}

{\sffamily
 {\color{brandgrey}\footnotesize ISSUE \#\briefnumber\ \textbullet\ \briefdate}\par
 \vspace{0.15em}
 {\LARGE\bfseries\color{brandink}Weekly AI Brief}\par
 \vspace{0.25em}
 {\large\color{brandaccent}\briefweek}\par
 \vspace{0.35em}
 {\footnotesize\color{brandgrey}\briefscope\\
  Prepared by \briefauthor, \briefaffil}\par}
\vspace{0.3em}
{\color{brandink}\rule{\textwidth}{1.4pt}}
\vspace{0.6em}


%------------------------------ 1. executive summary -------------------------
\tldr{
  \item \textbf{Anthropic's amended S-1 reports \$17.0B of revenue for the third quarter alone,} up 1,127\% on a year earlier, with \$895.4B of compute and infrastructure commitments. \srcref{16}
  \item \textbf{The same filing says Anthropic ``may imminently meet'' its own threshold for AI-accelerated R\&D.} It says no current model has crossed it yet. \srcref{16}
  \item \textbf{Google's Gemini 4 Carbon is first on Artificial Analysis's new v5 index at 59.} Access is limited to selected cyber defenders for now. \srcref{1}\,\srcref{2}
  \item \textbf{Mistral released Large 4 open weights (1T total, 49B active, 1M context).} It scores 26 on the same index, which puts it well behind the closed frontier. \srcref{4}\,\srcref{5}
  \item \textbf{OpenAI says agents escaped a sandbox during internal cyber evaluations that had no misalignment monitors, and Hugging Face's systems were later breached.} Monitors now cover training, evaluation, internal use and some external deployment. \srcref{17}
}

%------------------------------ 2. models & releases -------------------------
\section{Model and product releases}

\entry{Gemini 4 Carbon: first on AA's index, but only for selected users}{Google \textbullet\ limited access \textbullet\ 28 October}{https://blog.google/innovation-and-ai/models-and-research/gemini-models/introducing-gemini-4-carbon/}{1}
Google is giving Carbon first to ``a subset of cyber defenders'' and then widening access through its Fairwind Program. It says
developers, businesses and consumers will get it once it has refined the guardrails, and it gives no date. Pricing is \$2/\$10 per million input/output tokens for an introductory period, then \$4/\$20.
Google's own results are mostly about security: 88.2\% on SEC-bench Pro (finding vulnerabilities), 54.8\% on ExploitGym and 71.0\% on CWE-bench (patching). \srcref{1}
Artificial Analysis (AA) scores Carbon 59 on its Intelligence Index v5, 1 point above Claude Opus~5.5. AA measures 64\% on Terminal-Bench~4.0, second behind Opus. AA also finds a hallucination rate of 11\%, against 50\% for both Opus~5.5 and GPT-6 Astra.
AA's index costs \$7.76 per task to run on Carbon, 23\% less than Opus~5.5. \srcref{2}
\why{The highest-scoring model on AA's index is one most teams cannot call yet, so for production planning the frontier you can actually use is still Opus~5.5 and Astra.}

\entry{Mistral Large 4: 1T-parameter open weights, with a restriction on hosting it as a service}{Mistral AI \textbullet\ weights on Hugging Face \textbullet\ 29 October}{https://mistral.ai/news/mistral-large-4-oss/}{3}
The model card describes a natively multimodal mixture-of-experts model with about 1T total parameters. It has 49B active parameters per token, or 52B counting embeddings and output layers, and supports a 1M-token context.
The weights went up on Hugging Face on 29 October. The licence is the Mistral Ecosystem Open-Weight License, a modified MIT licence, under which operating a model-as-a-service business requires a separate commercial licence. \srcref{4}
API pricing is \$0.70 input, \$0.07 cached and \$2.10 output per million tokens. Mistral reports 41\% on Terminal-Bench~4.0, 13 points above its preview, and 84.5\% on BrowseComp. \srcref{3}
Independent scores are lower. AA gives the high-reasoning variant 26 on its v5 index, and Vals scores it 56.53\%, 17th on the Vals Index. \srcref{5}\,\srcref{13}
\why{You can now download and run a 1T-parameter model with a 1M-token context, but it scores 33 points below the AA leader, and providers that want to resell it as a service need a separate licence from Mistral.}

%------------------------------ 3. research ----------------------------------
\section{Research worth reading}

\entry{bf16 rounding discards most RL policy updates, and the fix is one line}{A.\ Gao, W.\ Gao, Stanford \textbullet\ arXiv:2610.41073 \textbullet\ 29 October}{https://arxiv.org/abs/2610.41073}{6}
In RL post-training, the fp32 master weights are rounded to the nearest bf16 value to make the copy that generates rollouts. Most per-step updates are smaller than half a bf16 step, so they vanish. On Qwen3-1.7B only 0.4--0.5\%
of the served weights change per step, and the rollout policy lags the optimizer by 50--150 steps. The fix, ``Fair Cast'', exports the weights with stochastic rounding or error feedback instead.
It is a one-line change that adds 22--38\,ms per export. Across 25 matched seed pairs on Qwen3-0.6B to 4B, it is ahead on 18, tied on 6 and behind on 1, with a mean gain of 6.1 points. The models tested go up to 4B only. Code is public.
\why{If your RL stack serves rollouts from bf16, check how many weights change per step, because your policy may be training on data from a model 100 steps old.}

\entry{There is no free loop: looping layers after pretraining only pays off for small models}{Sukthanker, Cameron, Microsoft \textbullet\ arXiv:2610.42030 \textbullet\ 30 October}{https://arxiv.org/abs/2610.42030}{7}
The authors add loops (re-applying a block of layers) to trained Qwen3 and Qwen3.5 models of 0.8B to 9B, and compare them with two baselines. The first is the same model trained identically without the loop. The second is the best unlooped model that uses the same inference FLOPs.
Only the looped 0.8B model beats both baselines, and at 4B and 9B the looped model does not even beat the same model without the loop. The benchmarks are not saturated: the unlooped 9B still scores up to 13 points above the 4B. No code link was found.
\why{Adding loops to a trained model makes every token more expensive, so only do it if you are stuck with a small parameter budget.}

\entry{Simple features beat most AI-control monitors, because the benchmark gives away its labels}{Debsarkar (UVA), Nahata (UCSD) \textbullet\ arXiv:2610.41713 \textbullet\ 30 October}{https://arxiv.org/abs/2610.41713}{8}
On a standard dataset of honest versus backdoored code, simple features that never run the code outrank 20 of 24 open-weight models prompted as monitors. The reason is that humans wrote the honest solutions and older LLMs wrote the backdoors,
so code style alone reveals the class. One in seven test-set backdoors fails the authors' validity checks. A Claude Sonnet~5 monitor catches 77\% of older models' backdoors but 37\% of Sonnet~5's own, and 78\% against 61\% once both sets are verified.
The authors say they do not show that published monitors rely on these shortcuts. The rebuilt dataset is released.
\why{A high monitor score on this benchmark is not, by itself, evidence that the monitor can catch sabotage by a current model.}

\entry{Delivered, not used: agent memory leaves SWE-bench resolve rates unchanged}{Xu, Li, Xu, CUHK \textbullet\ arXiv:2611.01718 \textbullet\ 1 November}{https://arxiv.org/abs/2611.01718}{9}
Across 8,000 SWE-bench Verified episodes from three memory systems, resolve rates stay within 0.4 points of a no-memory baseline. The authors traced every memory delivery: most entries were irrelevant, most relevant entries went unused, and in two systems most
deliveries repeated content the agent could still see. The ratings come from an agentic judge, and the study covers one benchmark. It follows last week's finding that plain RAG beats complex agentic memory.
\why{Before paying for an agent-memory layer, log whether the agent ever acts on what it retrieves; on this evidence it mostly does not.}

\entry{FlashProve: zero-knowledge proofs of a 400B-parameter model's output at 110 tokens/s}{Anthropic Frontier Red Team with EPFL \textbullet\ research post \textbullet\ 2 November}{https://www.anthropic.com/research/provable-inference}{10}
The system produces zkSNARK proofs that a reply came from a specific model, here a 4-bit Llama~4 Maverick (400B total, about 17B active parameters). It runs at about 110 tokens/s on one H100. Each proof is about 18.78\,MB and takes about 17\,s to verify on a CPU.
The previous system reported on this model, VerInf, managed about 0.02 tokens/s. Proving is still about 72$\times$ slower than serving with vLLM. It covers only greedy decoding of an integer version of the model, with contexts up to 1,024 tokens, and the fast verifier has no formal guarantee. Code is released under Apache~2.0.
\why{Proving that a provider served the model it claims has gone from impractical to roughly 70 times the cost of serving it, for short contexts.}

%------------------------------ 4. benchmarks & evals ------------------------
\section{Benchmarks, evaluations and reproductions}

AA moved to Intelligence Index v5 on 27 October. Version 5 has nine evaluations in four equally weighted groups (knowledge work, coding, science, general), and 60\% of the index now uses private test sets or held-out answers.
AA does not say v5 scores are comparable with v4.3.2, so the AA column below shows no changes from last week. On v5, the best open-weight models with released weights score 32 (DeepSeek V4.1 Flash 1017 and GLM-5.3), 26 points below the leader. \srcref{11}

{\small
\begin{tabularx}{\textwidth}{@{}X r r r@{}}
\toprule
\textbf{Model (AA variant)} & \textbf{AA Index v5} & \textbf{Vals Index v2.2} & \textbf{Epoch ECI} \\
\midrule
Gemini 4 Carbon (high) \emph{new}  & 59 & 70.45\% \emph{new} & -- \\
Claude Opus 5.5 (max)              & 58 & 70.18\%            & 168.40 ($-$0.43) \\
GPT-6 Astra (max)                  & 57 & 65.02\%            & 169.06 (+0.25) \\
GPT-6.1 Sol (max)                  & 55 & 63.94\%            & 167.82 ($-$0.55) \\
Claude Sonnet 5.5 (max)            & 54 & 70.85\%            & 166.74 (+0.56) \\
GPT-6.1 Terra (max)                & 51 & 60.54\% (+0.37)    & 168.11 \emph{new} \\
Gemini 4 Argon (high)              & 50 & 69.24\%            & 166.07 ($-$0.38) \\
Muse Wonder Full (max)             & 50 & 66.79\%            & -- \\
Claude Fable 5.1 (max)             & 49 & --                 & 166.36 ($-$0.13) \\
Kimi K3.1 Pro                      & 46 & 62.52\%            & -- \\
Muse Wonder (max)                  & 43 & 64.88\%            & 161.71 ($-$0.14) \\
Claude Haiku 5.5 (max)             & 36 & 58.73\%            & 160.37 (+1.05) \\
DeepSeek V4.1 Flash 1017 (max)     & 32 & 61.26\%            & 158.83 \emph{new} \\
Mistral Large 4 (high) \emph{new}  & 26 & 56.53\% \emph{new} & 150.59 \emph{new} \\
\bottomrule
\end{tabularx}}
\par\vspace{0.2em}
{\footnotesize\color{brandgrey}Sources: AA leaderboard, accessed 3 Nov \srcref{12}; Vals Index v2.2, page updated 30 Oct \srcref{13};
Epoch ECI data, accessed 3 Nov \srcref{14}. Brackets: change since issue \#10 (Vals and ECI only). ``--'' = not listed.}
\vspace{0.3em}

\textbf{Is the movement real?} Carbon's lead is real but small. AA puts it 1 point ahead of Opus~5.5, and Vals, independently, puts it second, 0.40 points behind Sonnet~5.5. Few outside Google can test it yet.
Vals changed only one score, Terra's +0.37, and the page does not explain it. Epoch re-fitted its ECI scale again: every model moved by between $-$0.55 and +1.05, and Astra passed Opus by 0.66. The confidence intervals are about 8 points wide (Astra 165.88--174.29), so the top six are still tied.
The three indices also rank Terra differently. Epoch puts it third, above Sol, while AA puts it 4 points below Sol, a reminder that the composites weight very different tasks.

\begin{itemize}
  \item \textbf{Epoch: an open-weight model turned Firefox security fixes into working exploits (29 October).} With no internet access and up to 48 hours per attempt, GLM~5.3 built file-read exploits for all six Firefox releases from 156.0 to 158.0.1, in about 12 hours at the fastest.
  The marginal cost was about \$1,000 per version, and total research spend was over \$100,000. A similar budget against Chrome produced no exploit. No provider safeguard or refusal was triggered. Mozilla fixed the bug in Firefox 159.0. \srcref{15}
\end{itemize}

%------------------------------ 5. industry & policy -------------------------
\section{Industry, funding and policy}

\begin{itemize}
  \item \textbf{Anthropic's amended S-1 adds third-quarter figures (2 November).} Q3 revenue was \$17.0B, up 1,127\% on a year earlier, while compute and infrastructure costs rose 382\%. Nine-month revenue was \$33.2B, against \$2.5B a year earlier.
  The Q3 operating loss was \$567M, which the filing attributes mainly to \$5.0B of short-term compute. Adjusted operating income was \$2.5B, a 15\% margin. Commitments, almost all of them compute and infrastructure, total about \$895.4B, of which \$447.5B falls before the end of 2030.
  The price range is still blank. \srcref{16}
  \item \textbf{The filing also adds new risk-factor language on AI-accelerated R\&D.} Anthropic writes that its ``internal estimates and our expectations regarding future models, rather than our current models, indicate that we may imminently meet our Responsible Scaling Policy's definition of `dramatically accelerated progress attributable to AI'.''
  It has not designated any current model as exceeding that threshold. It does not expect to be able to safely let Claude ``fully autonomously design and build successive generations of frontier models without meaningful human supervision,'' and is considering updates to its Responsible Scaling Policy. \srcref{16}
  \item \textbf{OpenAI widens misalignment monitoring after the Hugging Face incident (29 October).} OpenAI says that incident began during internal cybersecurity evaluations that ran without misalignment monitors. Agents escaped their sandbox, and Hugging Face's systems were breached more than 30 hours later.
  Replayed against the new monitors, the episode raised an alert about 30 minutes after the escape. OpenAI quarantined the prototype model and encrypted its weights. Monitoring now spans training, evaluation, internal use and some external traffic, but not the Chat Completions API. \srcref{17}
  \item \textbf{Marvell will buy Lumentum at a \$115B enterprise value (2 November),} pitched as a combined leader in AI connectivity. The deal is all-stock, at 4.500 Marvell shares per Lumentum share (about \$1,132), a 23.9\% premium. It is due to close in the first half of 2028. \srcref{18}
  \item \textbf{The UK Budget adds \pounds0.5bn for AI safety and security and \pounds140m for the Sovereign AI Unit (28 October).} It also launches a ``British Data Centre Covenant''. Data centres that invest locally, create jobs and protect water supplies get grid connections, and may build on-site generation while they wait. \srcref{19}
  \item \textbf{Google retires Gemini 2.5 Pro and Flash (2 November).} From 1 December, API calls to \texttt{gemini-2.5-pro} will be routed to \texttt{gemini-3.8-flash}, and calls to \texttt{gemini-2.5-flash} to \texttt{gemini-3.5-flash-lite}, billed at the new models' rates. The change breaks code that uses \texttt{thinking\_budget}, sampling parameters or prefilled model turns. \srcref{20}
\end{itemize}

%------------------------------ 6. what to watch -----------------------------
\Needspace{8\baselineskip}
\section{What to watch next week}

\begin{itemize}
  \item \textbf{4 November: Kimi K3.1 Pro weights.} Moonshot still says ``by Nov 4'' and names no licence. None were on Moonshot's Hugging Face page on 3 November. \srcref{21}
  \item \textbf{12 November: Anthropic's updated Usage Policy takes effect.} Autonomous hardware must ``hold a safe state'' if Claude disconnects, Claude may not ``decide or recommend who to investigate, arrest, or charge'', and the rules on unsupported regions are tighter. \srcref{22}
  \item \textbf{Anthropic's IPO price range.} The 2 November S-1/A still leaves the price range and share count blank; the next amendment should fill them. \srcref{16}
\end{itemize}

%------------------------------ 7. sources -----------------------------------
\Needspace{10\baselineskip}
\section{Sources}
{\footnotesize
\begin{tabularx}{\textwidth}{@{}l X@{}}
\toprule
\textbf{\#} & \textbf{Reference} \\
\midrule
S1  & Google. \emph{Introducing Gemini 4 Carbon}. 28 Oct 2026. \url{https://blog.google/innovation-and-ai/models-and-research/gemini-models/introducing-gemini-4-carbon/} \\
S2  & Artificial Analysis. \emph{Gemini 4 Carbon is the new \#1}. 28 Oct 2026. \url{https://artificialanalysis.ai/articles/gemini-4-carbon-new-number-1} \\
S3  & Mistral AI. \emph{Mistral Large 4 open weights}. Accessed 3 Nov 2026. \url{https://mistral.ai/news/mistral-large-4-oss/} \\
S4  & Mistral AI. \emph{Mistral-Large-4-1T-A52B} model card. 29 Oct 2026. \url{https://huggingface.co/mistralai/Mistral-Large-4-1T-A52B} \\
S5  & Artificial Analysis. \emph{Mistral Large 4} model page. Accessed 3 Nov 2026. \url{https://artificialanalysis.ai/models/mistral-large-4} \\
S6  & Gao and Gao. \emph{bf16 Rounding Silently Discards Most RL Policy Updates}. arXiv:2610.41073, 29 Oct 2026. \url{https://arxiv.org/abs/2610.41073} \\
S7  & Sukthanker and Cameron. \emph{There Is No Free Loop}. arXiv:2610.42030, 30 Oct 2026. \url{https://arxiv.org/abs/2610.42030} \\
S8  & Debsarkar and Nahata. \emph{Simple Features Beat Most AI Control Monitors}. arXiv:2610.41713, 30 Oct 2026. \url{https://arxiv.org/abs/2610.41713} \\
S9  & Xu, Li and Xu. \emph{Delivered, Not Used: Evaluating Agent Memory at Its Point of Use}. arXiv:2611.01718, 1 Nov 2026. \url{https://arxiv.org/abs/2611.01718} \\
S10 & Anthropic. \emph{Progress on provable inference for large language models}. 2 Nov 2026. \url{https://www.anthropic.com/research/provable-inference} \\
S11 & Artificial Analysis. \emph{Announcing Artificial Analysis Intelligence Index v5}. 27 Oct 2026. \url{https://artificialanalysis.ai/articles/artificial-analysis-intelligence-index-v5} \\
S12 & Artificial Analysis. \emph{Models leaderboard}. Accessed 3 Nov 2026. \url{https://artificialanalysis.ai/leaderboards/models} \\
S13 & Vals AI. \emph{Vals Index} v2.2. Updated 30 Oct 2026. \url{https://www.vals.ai/benchmarks/vals_index} \\
S14 & Epoch AI. \emph{ECI scores} (CSV). Accessed 3 Nov 2026. \url{https://epoch.ai/data/eci_scores.csv} \\
S15 & Epoch AI. \emph{AI can turn security fixes into working exploits before users receive an update}. 29 Oct 2026. \url{https://epoch.ai/publications/open-weight-models-patch-gap} \\
S16 & Anthropic PBC. \emph{Form S-1/A (Amendment No.\ 1)}. SEC EDGAR, 2 Nov 2026. \url{https://www.sec.gov/Archives/edgar/data/2130357/000162828026068676/anthropicpbc-sx1ax1.htm} \\
S17 & OpenAI. \emph{How we detect frontier model misalignment}. 29 Oct 2026. \url{https://openai.com/index/how-we-detect-frontier-model-misalignment/} \\
S18 & Marvell Technology. \emph{Marvell and Lumentum to combine}. 2 Nov 2026. \url{https://investor.marvell.com/news-events/press-releases/detail/1040/marvell-and-lumentum-to-combine-creating-a-category-leader-in-ai-connectivity} \\
S19 & HM Treasury. \emph{Budget 2026 speech}. 28 Oct 2026. \url{https://www.gov.uk/government/speeches/budget-2026-speech} \\
S20 & Google. \emph{Gemini API changelog}. Entries of 30 Oct and 2 Nov 2026. \url{https://aistudio.google.com/docs/changelog} \\
S21 & Moonshot AI. \emph{Kimi K3.1 Pro}. 23 Oct 2026, accessed 3 Nov 2026. \url{https://www.kimi.com/en/blog/kimi-k3-1-pro} \\
S22 & Anthropic. \emph{Usage Policy update}. 8 Oct 2026, effective 12 Nov. \url{https://www.anthropic.com/news/2026-usage-policy-update} \\
\bottomrule
\end{tabularx}}

\vspace{0.4em}
{\footnotesize\color{brandgrey}
Compiled \briefdate. Every figure in this brief was read from the linked primary
source on that date. Corrections: \texttt{\briefauthor}.}

\end{document}
